\chapter{几何专题}

\section{向量}
数学是从认识和研究图形和数开始的。大体上可以说，图形的优点是直观形象，能更直接地用来描述我们周围的世界，也更容易理解。但图形不便于用于计算，利用几何推理的方法来研究图形，灵活性、偶然性太大，不容易掌握。数的优点是有比较死板的方法进行运算，便于掌握，但比图形更抽象，将客观世界用数来描述的难度更大一些。笛卡尔引进了坐标之后，打破了数与形的界限，将几何图形最基本的元素——点用坐标来表示，将曲线、曲面用方程来表示，通过对坐标和方程的代数运算来研究几何图形的性质，这就是解析几何。

向量也是一种图形，具有直观形象的优点，便于用来描述客观世界。向量又可以直接进行运算。因此向量兼有图形和数的优点。我们可以将几何图形用向量来描述，通过向量的运算来解决几何问题；也可以进一步将向量的运算转化为坐标的运算来解决问题，这是建立解析几何的理论与算法体系的一种简明有效的途径。

向量是一种“价廉物美”的数学工具，学起来很容易，需要花费的时间和精力不多；用处又很大，可以解决很多的几何和代数的问题。

\subsection{向量的定义}
向量，一称矢量，是一类有大小和方向的量。每个向量\(\vec{a}\)可以用一条有向线段\(A B\)来表示，记作\(\vec{a} = \vec{A B}\)。从\(A\)到\(B\)的方向表示向量\(\vec{a}\)的方向；\(A B\)的长度\(\left| A B \right|\)表示向量\(\vec{a}\)的大小，称为\(\vec{a}\)的模，记作\(\left| \vec{a} \right|\)。当\(\left| \vec{a} \right| = \left| A B \right| = 0\)时\(A\)与\(B\)重合，\(\vec{a} = \vec{A A}\)，称为零向量，记为\(\vec{0}\)。零向量的方向不确定，可以为任意方向。

同一个向量\(\vec{a}\)可以由不同的有向线段\(A B\)、\(C D\)表示，只要\(A B\)、\(C D\)的方向相同，长度也相等，即使出发点\(A\)、\(C\)不同也有\(\vec{A B} = \vec{C D}\)。但从同一点出发表示同一个向量的有向线段是唯一确定的：\(\vec{O A} = \vec{O B} \Longleftrightarrow A\)与\(B\)重合。

当非零向量\(\vec{a}\)与\(\vec{b}\)的方向相同或相反时，表示它们的有向线段既可以共线，也可以平行，此时我们称\(\vec{a} , \vec{b}\)平行（也可以称\(\vec{a} , \vec{b}\)共线）。我们规定：零向量与任意向量平行。

\subsection{向量的加法}
设\(\vec{a} , \vec{b}\)是任意两个向量。从任意一点\(O\)出发，作\(\vec{O A} = \vec{a}\)，再从\(O A\)的终点\(A\)出发作\(\vec{A B} = \vec{b}\)，则定义\(\vec{a}\)与\(\vec{b}\)相加所得的和\(\vec{a} + \vec{b} = \vec{O A} + \vec{A B} = \vec{O B}\)。向量加法的这个法则称为三角形法则。如果将\(\vec{O A} , \vec{A B}\)理解为位移，则\(\vec{O A} + \vec{A B}\)就表示两个位移的合成：\(O \to A \to B\)总的效果是由\(O\)到\(B\)，也就是\(\vec{O B}\)。

对任意3个向量\(\vec{a} , \vec{b} , \vec{c}\)作\(\vec{a} = \vec{O A} , \vec{b} = \vec{A B} , \vec{c} = \vec{B C}\)，则易见\(\left( \vec{a} + \vec{b} \right) + \vec{c} = \vec{O C} = \vec{a} + \left( \vec{b} + \vec{c} \right)\)，这也说明了向量加法的结合律成立。

设\(\vec{a} = \vec{O A}\)与\(\vec{b} = \vec{A B}\)，再作\(\vec{O C} = \vec{b}\)，连接\(C B\)。则当\(O A , A B\)不共线时，\(O A B C\)是平行四边形，从而\(\vec{C B} = \vec{O A} = \vec{a}\)，如右图所示。当\(O A , A B\)共线时，也易见\(\vec{C B} = \vec{O A}\)。因此在所有情况下都有\(\vec{b} + \vec{a} = \vec{O C} + \vec{C B} = \vec{O B} = \vec{a} + \vec{b}\)。也即向量加法交换律成立。

\begin{wrapfigure}[6]{r}{0.38\linewidth}
    \centering
    \includegraphics[width=0.84\linewidth]{image10.png}
\end{wrapfigure}

右图中的等式\(\vec{a} + \vec{b} = \vec{O A} + \vec{O C} = \vec{O B}\)给出了求向量\(\vec{a} , \vec{b}\)的和的另一个法则，称为平行四边形法则：从同一点\(O\)出发，分别作\(\vec{O A} = \vec{a} , \vec{O C} = \vec{b}\)，作平行四边形\(O A B C\)，则对角线\(\vec{O B} = \vec{a} + \vec{b}\)。

易见\(\vec{a} + \vec{0} = \vec{0} + \vec{a} = \vec{a}\)对任意向量\(\vec{a}\)成立。

对每个向量\(\vec{a} = \vec{A B}\)，存在唯一一个与\(\vec{a}\)大小相等、方向相反的向量\(\vec{B A}\)，记作\(- \vec{a}\)，也就是\(\vec{B A} = - \vec{A B}\)，称为\(\vec{a}\)的负向量。我们有\(\vec{a} + \left( - \vec{a} \right) = \vec{A B} + \vec{B A} = \vec{A A} = \vec{0} , \left( - \vec{a} \right) + \vec{a} = \vec{B A} + \vec{A B} = \vec{B B} = \vec{0}\)。反过来，\(\vec{a}\)也是\(- \vec{a}\)的负向量，即\(- \left( - \vec{a} \right) = \vec{a}\)。

向量的减法：已知向量\(\vec{a} , \vec{b}\)，定义它们的差\(\vec{a} - \vec{b} = \vec{a} + \left( - \vec{b} \right)\)，它满足条件\(\left( \vec{a} - \vec{b} \right) + \vec{b} = \vec{a}\)。从同一点\(O\)出发，分别作\(\vec{O A} = \vec{a} , \vec{O B} = \vec{b}\)，则\(\vec{a} - \vec{b} = \vec{B A}\)。

\subsection{向量与实数相乘}
任一向量\(\vec{a}\)与任一实数\(k\)可以相乘，得到一个向量，记作\(k \vec{a}\)，定义如下：

\(k \vec{a}\)的长度\(\left| k \vec{a} \right| = \left| k \right| \left| \vec{a} \right|\)。当\(k > 0\)时，\(k \vec{a}\)与\(\vec{a}\)的方向相同；当\(k < 0\)时，\(k \vec{a}\)与\(\vec{a}\)的方向相反；当\(k = 0\)时，\(k \vec{a} = \vec{0}\)是零向量，方向不确定。

由向量与实数相乘的法则知道：向量\(\vec{a} , \vec{b}\)平行\(\Leftrightarrow\)其中一个向量可以写成另一个向量与某个实数的乘积。
\newpage
如下运算律对任意向量\(\vec{a} , \vec{b}\)和实数\(k , s\)成立：

（1）\(1 \vec{a} = \vec{a} , - 1 \vec{a} = - \vec{a} , k \left( s \vec{a} \right) = \left( k s \right) \vec{a}\)；

（2）（对实数加法的分配律）\(\left( k + s \right) \vec{a} = k \vec{a} + s \vec{a}\)；

（3）（对向量加法的分配律）\(k \left( \vec{a} + \vec{b} \right) = k \vec{a} + k \vec{b}\)。

长度为\(1\)的向量称为单位向量。当\(\vec{a} \ne \vec{0}\)时，\(\frac{1}{\left| \vec{a} \right|} \vec{a}\)就是与\(\vec{a}\)方向相同的单位向量。

\begin{example}
（1）已知\(\triangle A B C\)及其两边长\(\left| A B \right| , \left| A C \right|\)。怎样用向量的语言描述：点\(P\)在\(\angle B A C\)的角平分线上。

（2）\(O\)是平面上一定点，\(A , B , C\)是平面上不共线的三个点，动点\(P\)满足\(\vec{O P} = \vec{O A} + \lambda \left( \frac{\vec{A B}}{\left| \vec{A B} \right|} + \frac{\vec{A C}}{\left| \vec{A C} \right|} \right) , \lambda \in \lbrack 0 , + \infty \rightparen\)，则\(P\)点的轨迹一定过\(\triangle A B C\)的

（A）外心 \quad （B）内心 \quad （C）重心 \quad （D）垂心
\end{example}

\subsection{向量的数量积}
任意两个向量\(\vec{a} , \vec{b}\)可以作数量积（也称内积），得到一个实数：

\(\vec{a} \cdot \vec{b} = \left| \vec{a} \right| \left| \vec{b} \right| \cos \left\langle \vec{a} , \vec{b} \right\rangle\)，其中\(\left\langle \vec{a} , \vec{b} \right\rangle\)表示\(\vec{a} , \vec{b}\)的夹角。

物理例子：设物体在恒力\(\vec{F}\)作用下沿直线前进，产生位移\(\vec{s}\)，则\(\vec{F}\)所作功\(W = \vec{F} \cdot \vec{s}\)。

有两个特殊情况值得特别注意：

（1）记\(\vec{a}^{2} = \vec{a} \cdot \vec{a}\)，由于向量\(\vec{a}\)与自身夹角为0，由\(\cos 0 = 1\)知\(\vec{a}^{2} = \left| \vec{a} \right|^{2}\)；

（2）当\(\vec{a} \perp \vec{b}\)时，\(\cos \left\langle \vec{a} , \vec{b} \right\rangle = \cos \frac{\pi}{2} = 0\)，因此\(\vec{a} \cdot \vec{b} = 0\)；我们规定零向量与所有向量垂直，这样就有\(\vec{a} \cdot \vec{b} = 0 \Longleftrightarrow \vec{a} \perp \vec{b}\)。

利用数量积可以表示任意向量\(\vec{a}\)的长度\(\left| \vec{a} \right|\)和任意两个非零向量\(\vec{a} , \vec{b}\)所成的角\(\theta\)：
\begin{equation}
    \left| \vec{a} \right| = \sqrt{\vec{a}^{2}} ; \cos \theta = \frac{\vec{a} \cdot \vec{b}}{\left| \vec{a} \right| \left| \vec{b} \right|} = \frac{\vec{a} \cdot \vec{b}}{\sqrt{\vec{a}^{2}} \sqrt{\vec{b}^{2}}}
\end{equation}

向量的数量积满足一些与数的乘法类似的运算律：

（1）\(\vec{a} \cdot \vec{b} = \vec{b} \cdot \vec{a}\)；

（2）\(\left( \vec{a} + \vec{b} \right) \cdot \vec{c} = \vec{a} \cdot \vec{c} + \vec{b} \cdot \vec{c}\)；

（3）\(\left( k \vec{a} \right) \cdot \vec{b} = k \left( \vec{a} \cdot \vec{b} \right)\)。
\newpage
\begin{example}
对任意向量\(\vec{a} , \vec{b}\)证明乘法公式\({\left( \vec{a} - \vec{b} \right)}^{2} = \vec{a}^{2} + \vec{b}^{2} - 2 \vec{a} \cdot \vec{b}\)，并对\(\vec{a} , \vec{b}\)不平行的情形解释这个公式的几何意义。
\end{example}

{\kaiti{}思考\quad 为什么这里向量的运算和数的运算可以混为一谈。}

\begin{example}
设\(A B C D\)是三维空间中的不重合的任意四个点。求证：\(A C \perp B D \Longleftrightarrow \left| A B \right|^{2} + \left| C D \right|^{2} = \left| B C \right|^{2} + \left| D A \right|^{2}\)。

\end{example}

\begin{example}
\(A B C D\)是四面体，\(\angle B A C = \angle A C D\)，\(\angle A B D = \angle B D C\)。求证：\(\left| A B \right| = \left| C D \right|\)。
\end{example}

\subsection{向量的坐标表示法}
一般地，在平面上任意选定两个相互垂直的单位向量\(\vec{\imath} , \vec{\jmath}\)，则平面内每个向量\(\vec{a} = \vec{O P}\)都可以分解为与\(\vec{\imath}\)平行的向量\(\vec{O D}\)和与\(\vec{\jmath}\)平行的向量\(\vec{D P}\)之和。\(\vec{O D} , \vec{D P}\)可以分别写成\(\vec{\imath} , \vec{\jmath}\)的实数倍\(x \vec{\imath} , y \vec{\jmath}\)。于是\(\vec{a} = \vec{O P}\)被分解成\(\vec{\imath} , \vec{\jmath}\)的实数倍之和的形式\(\vec{O P} = x \vec{\imath} + y \vec{\jmath}\)，其中的两个系数\(x , y\)由\(\vec{O P}\)唯一决定，可以认为是用\(\vec{\imath} , \vec{\jmath}\)这两个单位向量一起去“度量”\(\vec{O P}\)得到的“量数”。我们称\(\vec{\imath} , \vec{\jmath}\)组成平面上的一组基，\(\left( x , y \right)\)为\(\vec{a} = \vec{O P}\)在这组基下的坐标。坐标\(\left( x , y \right)\)唯一地代表了向量\(\vec{a}\)，我们可以直接写\(\vec{a} = \left( x , y \right)\)。将向量用坐标表示之后，向量的运算就可以转化为坐标的运算。对向量\(\vec{a} = \left( x_{1} , y_{1} \right) , \vec{b} = \left( x_{2} , y_{2} \right)\)及实数\(k\)，我们有
\begin{equation}
    \vec{a} \pm \vec{b} = \left( x_{1} \vec{\imath} + y_{1} \vec{\jmath} \right) \pm \left( x_{2} \vec{\imath} + y_{2} \vec{\jmath} \right) = \left( x_{1} \pm x_{2} \right) \vec{\imath} + \left( y_{1} \pm y_{2} \right) \vec{\jmath}
\end{equation}
\begin{equation}
    k \vec{a} = k \left( x \vec{\imath} + y \vec{\jmath} \right) = \left( k x \right) \vec{\imath} + \left( k y \right) \vec{\jmath}
\end{equation}
\begin{equation}
    \vec{a} \cdot \vec{b} = \left( x_{1} \vec{\imath} + y_{1} \vec{\jmath} \right) \cdot \left( x_{2} \vec{\imath} + y_{2} \vec{\jmath} \right) = x_{1} x_{2} + y_{1} y_{2}
\end{equation}
可以用坐标的形式简写这三个式子：
\begin{equation}
    \left( x_{1} , y_{1} \right) + \left( x_{2} , y_{2} \right) = \left( x_{1} + x_{2} , y_{1} + y_{2} \right) , k \left( x , y \right) = \left( k x , k y \right) , \left( x_{1} , y_{1} \right) \cdot \left( x_{2} , y_{2} \right) = x_{1} x_{2} + y_{1} y_{2}
\end{equation}

由向量的数量积公式立即得到：

若\(\vec{a} = \left( x , y \right)\)，则\(\left| \vec{a} \right| = \sqrt{x^{2} + y^{2}}\)；

若\(\vec{a} = \left( x_{1} , y_{1} \right) , \vec{b} = \left( x_{2} , y_{2} \right)\)，则\(\cos \left\langle \vec{a} , \vec{b} \right\rangle = \frac{x_{1} x_{2} + y_{1} y_{2}}{\sqrt{x_{1}^{2} + y_{1}^{2}} \sqrt{x_{2}^{2} + y_{2}^{2}}}\)；

向量\(\vec{a} = \left( x_{1} , y_{1} \right) , \vec{b} = \left( x_{2} , y_{2} \right)\)垂直的充分必要条件是\(x_{1} x_{2} + y_{1} y_{2} = 0\)。
\newpage
\begin{exercise}
    \mbox{}

    （1）设\(A , B , C\)是实数，且\(A , B\)不全为\(0\)，求方程\(A x + B y + C = 0\)的图像。

    （2）已知直线\(l_{1} : A_{1} x + B_{1} y + C_{1} = 0\)与\(l_{2} : A_{2} x + B_{2} y + C_{2} = 0\)。求\(l_{1} , l_{2}\)的夹角\(\theta\)以及\(l_{1} \perp l_{2}\)的充分必要条件。

    （3）已知直线\(l_{1} , l_{2}\)的斜率分别是\(k_{1} , k_{2}\)，求\(l_{1} , l_{2}\)的夹角\(\theta\)以及\(l_{1} \perp l_{2}\)的充分必要条件。

    （4）设\(\alpha\)是任意角，\(n\)是任意正整数且\(\cos \alpha + \cos 2 \alpha + \cdots + \cos n \alpha \ne 0\)。化简
    \begin{equation*}
        \frac{\sin \alpha + \sin 2 \alpha + \cdots + \sin n \alpha}{\cos \alpha + \cos 2 \alpha + \cdots + \cos n \alpha} .
    \end{equation*}
\end{exercise}

可以很自然地将二维平面的向量坐标表示扩展到三维空间上甚至更高维度上。以三维空间为例。在空间中任意取定3个两两垂直的单位向量\(\vec{\imath} , \vec{\jmath} , \vec{k}\)，则空间中每个向量\(\vec{a}\)可以唯一地写成\(\vec{a} = x \vec{\imath} + y \vec{\jmath} + z \vec{k}\)的形式，其中\(x , y , z\)是实数，\(\vec{\imath} , \vec{\jmath} , \vec{k}\)称为空间中的一组基，\(\left( x , y , z \right)\)称为\(\vec{a}\)在这组基下的坐标，可以看成是用\(\vec{\imath} , \vec{\jmath} , \vec{k}\)这一组单位向量去“度量”\(\vec{a}\)得到的“量数”。\(\vec{a}\)可以用它的坐标\(\left( x , y , z \right)\)完全代表，直接写为\(\vec{a} = \left( x , y , z \right)\)。与平面向量类似地有空间向量的坐标运算公式：
\begin{equation*}
    \left( x_{1} , y_{1} , z_{1} \right) + \left( x_{2} , y_{2} , z_{2} \right) = \left( x_{1} + x_{2} , y_{1} + y_{2} , z_{1} + z_{2} \right) , \quad k \left( x , y , z \right) = \left( k x , k y , k z \right)
\end{equation*}
\begin{equation*}
    \left( x_{1} , y_{1} , z_{1} \right) \cdot \left( x_{2} , y_{2} , z_{2} \right) = x_{1} x_{2} + y_{1} y_{2} + z_{1} z_{2} , \quad \vec{a} = \left( x , y , z \right) \to \left| \vec{a} \right| = \sqrt{x^{2} + y^{2} + z^{2}}
\end{equation*}

向量\(\vec{a} = \left( x_{1} , y_{1} , z_{1} \right) , \vec{b} = \left( x_{2} , y_{2} , z_{2} \right)\)垂直的充分必要条件是\(x_{1} x_{2} + y_{1} y_{2} + z_{1} z_{2} = 0\)。

\begin{exercise}
    \mbox{}

    \noindent\begin{minipage}[t]{0.58\linewidth}
    \setlength{\parindent}{2em}\setlength{\parskip}{6pt}
    已知正方体\(A B C D - A^{'} B^{'} C^{'} D^{'}\)如图，求：

    （1）直线\(A B\)与平面\(B D A^{'}\)所成的角；

    （2）平面\(A^{'} B^{'} C D\)与\(B D D^{'} B^{'}\)的夹角。
    \end{minipage}\hfill
    \begin{minipage}[t]{0.34\linewidth}
    \centering
    \adjustbox{valign=t}{\includegraphics[width=0.92\linewidth]{image11.png}}
    \end{minipage}
\end{exercise}

\newpage

\section{三维空间中的几何关系}
三维空间中的几何关系主要包含点、线、面之间的位置关系，主要涉及距离、夹角等度量关系。为叙述方便，以下几何关系将在笛卡尔直角坐标系下使用向量语言进行叙述。对于其他坐标系下的描述这里暂不涉及，可由坐标变换关系自行推导。

\subsection{点与点的关系}
点与点的关系较为简单，一般考虑两点间的距离公式：若空间中两点\(A , B\)坐标分别为\(A \left( x_{1} , y_{1} , z_{1} \right) , B \left( x_{2} , y_{2} , z_{2} \right)\)，则\(\left| A B \right| = \left| \vec{A B} \right| = \left| \left( x_{2} - x_{1} , y_{2} - y_{1} , z_{2} - z_{1} \right) \right| = \sqrt{{\left( x_{2} - x_{1} \right)}^{2} + {\left( y_{2} - y_{1} \right)}^{2} + {\left( z_{2} - z_{1} \right)}^{2}}\)。

\subsection{点与线的关系}
如果点\(P\)的坐标为\(\left( x_{0} , y_{0} , z_{0} \right)\)，直线\(L\)的方向向量为\(\vec{l} = \left( a , b , c \right)\)，且直线上有一点\(Q\)的坐标为\(\left( x_{1} , y_{1} , z_{1} \right)\)，则点\(P\)到直线\(L\)的距离\(d\)为：
\begin{equation}
    d = \frac{\left| \vec{P Q} \times \vec{l} \right|}{\left| \vec{l} \right|}
\end{equation}

其中\(\vec{P Q} = \left( x_{1} - x_{0} , y_{1} - y_{0} , z_{1} - z_{0} \right)\)。定义两个三维向量\(\vec{a} , \vec{b}\)的外积\(\vec{a} \times \vec{b}\)（一称向量积）：两个三维向量的外积为一个向量，其模长为\(\left| \vec{a} \right| \left| \vec{b} \right| \sin \left\langle \vec{a} , \vec{b} \right\rangle\)，方向垂直于\(\vec{a} , \vec{b}\)，按从\(\vec{a}\)到\(\vec{b}\)的右手定则确定方向（当\(\vec{a} \parallel \vec{b}\)时，规定\(\vec{a} \times \vec{b} = \vec{0}\)）。

\subsection{点与面的关系}
如果点\(P\)的坐标为\(\left( x_{0} , y_{0} , z_{0} \right)\)，平面\(\pi\)的方程为\(A x + B y + C z + D = 0\)（\(A , B , C\)不全为0），则\(P\)到平面\(\pi\)的距离\(d\)为：
\begin{equation}
    d = \frac{\left| A x_{0} + B y_{0} + C z_{0} + D \right|}{\sqrt{A^{2} + B^{2} + C^{2}}}
\end{equation}

\subsection{线与线的关系}
如果两直线的方向向量分别为\(\vec{d}_{1} = \left( a_{1} , b_{1} , c_{1} \right)\)和\(\vec{d}_{2} = \left( a_{2} , b_{2} , c_{2} \right)\)，则两直线夹角\(\theta\)的余弦值为：
\begin{equation}
    \cos \theta = \frac{\vec{d}_{1} \cdot \vec{d}_{2}}{\left| \vec{d}_{1} \right| \left| \vec{d}_{2} \right|}
\end{equation}

\subsection{线与面的关系}
如果直线的方向向量为\(\vec{d} = \left( a , b , c \right)\)，平面的法向量为\(\vec{n} = \left( A , B , C \right)\)，则直线与平面的夹角\(\theta\)的正弦值为：
\begin{equation}
    \sin \theta = \frac{\vec{d} \cdot \vec{n}}{\left| \vec{d} \right| \left| \vec{n} \right|}
\end{equation}

\subsection{面与面的关系}
如果两平面的法向量分别为\(\vec{n}_{1} = \left( a_{1} , b_{1} , c_{1} \right)\)和\(\vec{n}_{2} = \left( a_{2} , b_{2} , c_{2} \right)\)，则两平面夹角\(\theta\)的余弦值为：
\begin{equation}
    \cos \theta = \frac{\vec{n}_{1} \cdot \vec{n}_{2}}{\left| \vec{n}_{1} \right| \left| \vec{n}_{2} \right|}
\end{equation}

于是可以容易地将三维空间中的线、面之间的平行、垂直等几何关系使用向量语言进行描述：

面与面平行/垂直：其法向量平行/垂直。

线与线平行/垂直：其方向向量平行/垂直。

线与面平行/垂直：线的方向向量与面的法向量垂直/平行。

因此，常见立体几何的题目的一般做法为：

第一步：建立合适的坐标系，使得题目所涉及各点坐标尽可能简单。

第二步：将待求几何量（一般为角度）或待证明几何关系用向量语言进行描述。

第三步：求相关向量的坐标表示。

（a）直线的方向向量：直接选取直线上已知的两点坐标计算；

（b）平面的法向量：

1）设其法向量为\(\vec{n} = \left( x , y , z \right)\)，利用平面已知的三点构建在该平面内的两个不平行的向量\(\vec{\alpha} , \vec{\beta}\)（或其他点，总的原则是选取使向量的坐标表示尽可能简单的点）。

2）列式计算\(\vec{n} \cdot \vec{\alpha} = 0\)和\(\vec{n} \cdot \vec{\beta} = 0\)，令\(x\)（或\(y , z\)）为一个合适的值（总的原则是选取容易求解方程的值，或者使\(\vec{n}\)成为单位向量），求出该平面的一个法向量。

第四步：将第三步得到的向量坐标代入第二步的向量语言描述中，进行相关的运算或列方程/等式，完成证明。

\section{线性空间*}
在讨论向量时，我们为向量定义了加法与数乘，并逐一验证了它们满足的运算律。在第二章中我们看到，函数之间可以逐点相加，一个函数也可以逐点与实数相乘：
\begin{equation}
    \left( f + g \right) \left( x \right) = f \left( x \right) + g \left( x \right) , \quad \left( k f \right) \left( x \right) = k f \left( x \right) .
\end{equation}

向量是有大小和方向的量，函数是量与量之间的对应关系。两者的运算规律一致：加法都满足交换律、结合律，都有零向量与负向量；数乘都满足结合律与分配律，等等。既然规律一致，凡是只用到这些规律的推导，对两类对象就都成立。把共同的运算规律提炼出来，作为一类对象的共同规定，就得到线性空间。

\subsection{线性空间的定义}

\begin{definition}\label{def:linspace}
设\(V\)是一个集合，在\(V\)上定义了加法与数乘两种运算：\(V\)中任意两个元素\(a , b\)对应\(V\)中一个元素\(a + b\)；任意实数\(k\)与任意元素\(a\)对应\(V\)中一个元素\(k a\)。如果对任意\(a , b , c \in V\)和任意实数\(k , s\)，都满足下列八条运算律：

（1）\(a + b = b + a\)；

（2）\(\left( a + b \right) + c = a + \left( b + c \right)\)；

（3）存在零向量\(\vec{0} \in V\)，使对一切\(a \in V\)都有\(a + \vec{0} = a\)；

（4）对每个\(a \in V\)，存在负向量\(- a \in V\)，使\(a + \left( - a \right) = \vec{0}\)；

（5）\(1 a = a\)；

（6）\(k \left( s a \right) = \left( k s \right) a\)；

（7）\(k \left( a + b \right) = k a + k b\)；

（8）\(\left( k + s \right) a = k a + s a\)。\\
那么就称\(V\)为一个线性空间（也称向量空间），\(V\)中的元素统称为向量。
\end{definition}
这八条运算律没有提出新的要求，它们就是向量与函数已经满足的规律。变化的是“向量”一词的含义：从现在起，线性空间中的元素都称为向量，不论它是实数、有向线段、数组还是函数。统一命名的好处是，只要一个命题是直接从这八条运算律推出的，它便对一切线性空间成立，不必对每个例子重新论证一遍。

{\kaiti{}思考\quad 只用八条运算律证明：零向量是唯一的；每个向量的负向量也是唯一的。\\ \hspace*{\parindent}再证明：对任意向量\(a\)和实数\(k\)，都有\(0 a = \vec{0}\)，\(\left( - 1 \right) a = - a\)。}

下面给出两个例子。

（1）\(\mathbb{R}^{n}\)。全体\(n\)元实数组\(\left( x_{1} , x_{2} , … , x_{n} \right)\)组成的集合记作\(\mathbb{R}^{n}\)，其中加法与数乘都按分量规定：
\begin{equation}
    \left( x_{1} , … , x_{n} \right) + \left( y_{1} , … , y_{n} \right) = \left( x_{1} + y_{1} , … , x_{n} + y_{n} \right) , \quad k \left( x_{1} , … , x_{n} \right) = \left( k x_{1} , … , k x_{n} \right) .
\end{equation}
容易验证，八条运算律逐条成立：每一条不过是把实数的相应运算律在每个分量上重复了一遍。前面讲向量的坐标表示时写下的坐标运算公式，就是\(n = 2\)或\(3\)时这里的两个式子。

（2）函数空间。设\(A\)是实数集的一个子集，把定义在\(A\)上的一切实值函数放在一起，加法与数乘都按逐点方式规定。容易逐条验证，它也是线性空间。零向量是恒等于零的函数，函数\(f\)的负向量是\(- f\)，其余各条都化归为实数的运算律。特别地，数列是定义在正整数集合上的函数，因此全体数列也构成一个线性空间。

{\kaiti{}思考\quad 在正实数集上规定“加法”为乘法（\(u \oplus v = u v\)）、“数乘”为乘方（\(k \odot u = u^{k}\)），这仍满足八条运算律吗？如果满足，它的“零向量”是哪个数？“负向量”又是什么？}

\subsection{子空间}

线性空间内部有时还含有更小的线性空间。例如，把平面上全体向量看作线性空间，过原点的一条直线自身也构成一个线性空间：直线上任意两个向量的和仍在直线上，任意实数倍也仍在直线上。这样的子集称为子空间。

\begin{definition}\label{def:subspace}
设\(U\)是线性空间\(V\)的一个子集。如果\(U\)在\(V\)的加法与数乘下自身也构成一个线性空间，就称\(U\)是\(V\)的子空间。
\end{definition}

判断一个子集是不是子空间，不必把八条运算律重新检查一遍，只需验证三条：

（1）零向量在\(U\)中：\(\vec{0} \in U\)；

（2）对加法封闭：若\(a , b \in U\)，则\(a + b \in U\)；

（3）对数乘封闭：若\(a \in U\)而\(k\)为实数，则\(k a \in U\)。

事实上，只要这三条成立，其余各条就自动成立。八条运算律对\(V\)中一切元素成立，自然也对\(U\)中的元素成立。而负向量由\(\left( - 1 \right) a = - a\)与对数乘封闭性得到（见前面的思考）。第（1）条也可以换成“\(U\)非空”：只要\(U\)含某个元素\(a\)，就有\(0 a = \vec{0} \in U\)。

例如，在\(\mathbb{R}^{3}\)中，过原点的直线与过原点的平面都是子空间。而不过原点的直线、平面不是子空间，因为它们不含零向量。注意到，本章前面用过的平面方程\(A x + B y + C z + D = 0\)中，平面过原点，当且仅当\(D = 0\)。

再来看线性方程组的解集。每个方程都是一次的方程组称为{\heiti{}线性方程组}，一般写成
\begin{equation}
    \left\{ \begin{matrix}a_{11} x_{1} + a_{12} x_{2} + \cdots + a_{1 n} x_{n} = b_{1} , \\ a_{21} x_{1} + a_{22} x_{2} + \cdots + a_{2 n} x_{n} = b_{2} , \\ \qquad \cdots \cdots \cdots \\ a_{m 1} x_{1} + a_{m 2} x_{2} + \cdots + a_{m n} x_{n} = b_{m} ,\end{matrix} \right.
\end{equation}
其中\(x_{1} , … , x_{n}\)是未知量，\(a_{i j}\)是第\(i\)个方程中\(x_{j}\)的系数，\(b_{1} , … , b_{m}\)是常数项。常数项全为零的线性方程组称为{\heiti{}齐次线性方程组}。设\(u = \left( u_{1} , … , u_{n} \right)\)与\(v = \left( v_{1} , … , v_{n} \right)\)是同一个齐次线性方程组的两个解。取组中任意一个方程，例如第\(i\)个方程
\(a_{i 1} x_{1} + a_{i 2} x_{2} + \cdots + a_{i n} x_{n} = 0\)，
把\(u + v\)的坐标代入它的左边，得到
\(a_{i 1} \left( u_{1} + v_{1} \right) + \cdots + a_{i n} \left( u_{n} + v_{n} \right) = \left( a_{i 1} u_{1} + \cdots + a_{i n} u_{n} \right) + \left( a_{i 1} v_{1} + \cdots + a_{i n} v_{n} \right) = 0\)，
所以\(u + v\)也满足这个方程。同理，把\(u\)的坐标同乘\(k\)，可知\(k u\)也满足它。这对每一个方程都成立，故\(u + v\)与\(k u\)仍是方程组的解。又零向量显然是解。于是齐次线性方程组的解集是\(\mathbb{R}^{n}\)的一个子空间。相反，如果方程组不是齐次的（常数项不全为零），那么零向量不是解，解集就不再是子空间。而当它有解时，它的解集可以看成齐次方程组的解集整体平移一个特解的结果。这种“平移过的子空间”，在讲仿射变换时还会出现。

函数空间中也有子空间。奇函数的全体是一个子空间：两个奇函数之和仍是奇函数，奇函数的实数倍仍是奇函数。偶函数的全体同理。此外，任何线性空间都有两个平凡的子空间：只含零向量的集合\(\left\{ \vec{0} \right\}\)和它自身。

\subsection{线性组合与线性无关}

设\(a_{1} , a_{2} , … , a_{m}\)是线性空间\(V\)中的一组向量。把每个向量各乘上一个实数再求和，得到的向量
\(k_{1} a_{1} + k_{2} a_{2} + \cdots + k_{m} a_{m}\)
称为这组向量的一个线性组合，其中\(k_{1} , … , k_{m}\)为任意实数。一组向量的所有线性组合组成的集合，称为由这组向量张成的子空间，记作\(\operatorname{span}\left( a_{1} , … , a_{m} \right)\)。

张成的子空间是线性子空间：两个线性组合之和仍是线性组合，把对应系数相加即可。线性组合的实数倍仍是线性组合，各系数全取零时得到零向量。此外，它还是包含\(a_{1} , … , a_{m}\)的最小子空间。任何包含这组向量的子空间，出于封闭性，必然包含它们的一切线性组合。

例如，平面上两个不共线的向量张成整个平面，平面内每个向量都可以写成这两个向量的实数倍之和，这就是坐标表示一节的内容。而一个非零向量的所有线性组合是一条过原点的直线。空间中三个不共面的向量张成整个空间，两个不共线的向量张成一个平面子空间。

再举一个具体计算：在\(\mathbb{R}^{3}\)中取\(a_{1} = \left( 2 , 1 , - 3 \right)\)，\(a_{2} = \left( 1 , - 2 , 4 \right)\)，由于
\(\left( 17 , - 4 , 2 \right) = 6 a_{1} + 5 a_{2}\)，
所以\(\left( 17 , - 4 , 2 \right)\)属于\(a_{1} , a_{2}\)张成的子空间。

下面讨论表示是否唯一的问题。同一个向量\(x\)用\(a_{1} , … , a_{m}\)的线性组合来表示，如果存在两种表示
\(x = k_{1} a_{1} + \cdots + k_{m} a_{m} = s_{1} a_{1} + \cdots + s_{m} a_{m}\)，
两式相减就得到
\(\left( k_{1} - s_{1} \right) a_{1} + \cdots + \left( k_{m} - s_{m} \right) a_{m} = \vec{0}\)。
可见，表示是否唯一，取决于零向量能否用一组不全为零的系数表示出来。把这种情形单独命名：

\begin{definition}\label{def:linindep}
设\(a_{1} , … , a_{m}\)是线性空间\(V\)中的一组向量。如果只有当\(k_{1} = k_{2} = \cdots = k_{m} = 0\)时，才有
\(k_{1} a_{1} + k_{2} a_{2} + \cdots + k_{m} a_{m} = \vec{0}\)，
就称这组向量线性无关；否则称它们线性相关，也就是说，存在一组不全为零的实数\(k_{1} , … , k_{m}\)，使得上式成立。
\end{definition}

由定义可直接得到两条常用的结论。

其一，含零向量的向量组一定线性相关：取零向量的系数为\(1\)、其余系数为\(0\)，就得到一个系数不全为零的零组合。

其二，向量组线性相关，当且仅当其中至少有一个向量能由其余向量线性表示。事实上，若某个零组合中\(k_{j} \ne 0\)，把含\(a_{j}\)的项移到等式另一边再除以\(k_{j}\)，就得到\(a_{j}\)用其余向量表示的式子；反过来，若\(a_{j}\)能由其余向量线性表示，移项即得一个零组合，其中\(a_{j}\)的系数为\(1\)，不为零。

这两条说明：线性相关的向量组可以去掉某个向量，张成的子空间不变；线性无关的向量组去掉任何一个向量，张成的子空间都会变小。

在平面和空间中，线性相关与线性无关有直观的几何意义：两个向量线性相关，当且仅当它们共线（方向相同或相反）；三个向量线性相关，当且仅当它们共面。

例如，在\(\mathbb{R}^{3}\)中，向量组\(\left( 2 , 3 , 1 \right) , \left( 1 , - 1 , 2 \right) , \left( 7 , 3 , 8 \right)\)线性相关，因为
\(2 \left( 2 , 3 , 1 \right) + 3 \left( 1 , - 1 , 2 \right) - \left( 7 , 3 , 8 \right) = \vec{0}\)。
而\(\left( 1 , 0 , 0 \right) , \left( 0 , 1 , 0 \right) , \left( 0 , 0 , 1 \right)\)线性无关：若
\(k_{1} \left( 1 , 0 , 0 \right) + k_{2} \left( 0 , 1 , 0 \right) + k_{3} \left( 0 , 0 , 1 \right) = \vec{0}\)，
则\(\left( k_{1} , k_{2} , k_{3} \right) = \left( 0 , 0 , 0 \right)\)，即\(k_{1} = k_{2} = k_{3} = 0\)。

\subsection{基与坐标}

\begin{definition}\label{def:basis}
设\(e_{1} , … , e_{n}\)是线性空间\(V\)中的一组向量。如果这组向量线性无关，并且张成\(V\)（即\(V\)中每个向量都能表示为它们的线性组合），就称\(e_{1} , … , e_{n}\)为\(V\)的一组基。
\end{definition}

由前面的讨论，基可以用一句话重新描述：\(e_{1} , … , e_{n}\)是\(V\)的基，当且仅当\(V\)中的每个向量都能唯一地表示为它们的线性组合。下面说明这一点：若\(e_{1} , … , e_{n}\)线性无关，而同一个向量\(x\)有两种表示
\(x = k_{1} e_{1} + \cdots + k_{n} e_{n} = s_{1} e_{1} + \cdots + s_{n} e_{n}\)，
两式相减得\(\left( k_{1} - s_{1} \right) e_{1} + \cdots + \left( k_{n} - s_{n} \right) e_{n} = \vec{0}\)，故两个表示的所有系数都相同。反过来，若\(e_{1} , … , e_{n}\)线性相关，则存在不全为零的\(k_{1} , … , k_{n}\)使\(k_{1} e_{1} + \cdots + k_{n} e_{n} = \vec{0}\)。于是零向量既能写成系数全为零的组合，又能写成这组系数不全为零的组合，唯一性被破坏。

设\(e_{1} , … , e_{n}\)是\(V\)的一组基。对\(V\)中的每个向量\(x\)，唯一确定的系数\(\left( k_{1} , … , k_{n} \right)\)（使\(x = k_{1} e_{1} + \cdots + k_{n} e_{n}\)）称为\(x\)在这组基下的坐标。

在本章讲向量的坐标表示时，我们取定两个互相垂直的单位向量\(\vec{\imath} , \vec{\jmath}\)作为基，把每个平面向量写成\(x \vec{\imath} + y \vec{\jmath}\)，其中的系数就是坐标。现在看到，只要线性无关且张成整个空间，任何一组向量都可以充当基。注意到，同一个向量在不同基下的坐标一般不同，坐标是相对于基而言的。

例如，在\(\mathbb{R}^{2}\)中，\(e_{1} = \left( 1 , 1 \right)\)与\(e_{2} = \left( 1 , - 1 \right)\)线性无关，并且对任意\(\left( x , y \right)\)有
\(\left( x , y \right) = \frac{x + y}{2} e_{1} + \frac{x - y}{2} e_{2}\)，
所以它们构成\(\mathbb{R}^{2}\)的一组基。向量\(\left( 3 , 1 \right) = 2 e_{1} + e_{2}\)在这组基下的坐标是\(\left( 2 , 1 \right)\)。而在以\(\vec{\imath} , \vec{\jmath}\)为基的坐标下，同一个向量的坐标是\(\left( 3 , 1 \right)\)。此外，把向量写成坐标之后，向量的加法与数乘就转化为数组的加法与数乘，这与前面坐标表示一节中的做法一致。

\subsection{维数}

一组基里应当含有多少个向量？不同的基看起来可以很不一样：平面上既可以用两个互相垂直的单位向量作基，也可以用任意两个不共线的向量作基。但它们所含向量的个数总是一样的。下面证明这一点。

\begin{theorem}\label{thm:dim}
设\(e_{1} , … , e_{n}\)与\(f_{1} , … , f_{m}\)都是线性空间\(V\)的基，则\(n = m\)。
\end{theorem}

证明的关键是下面这个事实：若\(u_{1} , … , u_{p}\)线性无关，并且每个\(u_{i}\)都能由\(w_{1} , … , w_{q}\)线性表示，则\(p \le q\)。为说明它，我们把\(w\)组中的向量逐个换成\(u\)组中的向量，并始终保持当前向量组张成\(V\)这一性质。

第一步：把\(u_{1}\)换进去。因为\(u_{1}\)可由\(w_{1} , … , w_{q}\)线性表示，而\(u_{1} \ne \vec{0}\)（线性无关的向量组中不含零向量），所以表示式中至少有一个\(w\)的系数不为零。不妨设\(w_{1}\)的系数不为零，于是\(w_{1}\)也可由\(u_{1} , w_{2} , … , w_{q}\)线性表示。这说明：凡能由\(w_{1} , … , w_{q}\)线性表示的向量，把其中的\(w_{1}\)替换成它的表示式，就能由\(u_{1} , w_{2} , … , w_{q}\)线性表示。也就是说，用\(u_{1}\)替换\(w_{1}\)后，新向量组仍然张成\(V\)。

第\(k\)步：假设已经把\(u_{1} , … , u_{k - 1}\)换入，当前向量组仍张成\(V\)，且其中仍有\(w\)。那么\(u_{k}\)可由当前向量组线性表示（因为当前向量组张成\(V\)，而\(u_{k} \in V\)）。如果表示式中所有\(w\)的系数都为零，则\(u_{k}\)可由\(u_{1} , … , u_{k - 1}\)线性表示，这与\(u_{1} , … , u_{p}\)线性无关矛盾；所以表示式中至少有一个\(w\)的系数不为零，把该\(w\)换成\(u_{k}\)，由第一步同样的理由，新向量组仍张成\(V\)。

这样每换一步，\(u\)就多一个、\(w\)就少一个。如果\(p > q\)，那么进行\(q\)步之后，\(w\)已经全部被换出，剩下的\(u_{1} , … , u_{q}\)张成\(V\)，于是\(u_{q + 1}\)可由\(u_{1} , … , u_{q}\)线性表示，这与\(u_{1} , … , u_{p}\)线性无关矛盾。因此\(p \le q\)，事实得证。

现在把它用在两组基上：令\(u\)组为\(e_{1} , … , e_{n}\)，\(w\)组为\(f_{1} , … , f_{m}\)。\(e\)组线性无关，而\(f\)组张成\(V\)，所以\(e\)组中每个向量都能由\(f\)组线性表示，于是\(n \le m\)。再把两组的角色互换，又得\(m \le n\)。因此\(n = m\)：一组基所含向量的个数与基的选择无关，只由线性空间\(V\)本身决定。这个数称为\(V\)的维数，记作\(\dim V\)。（有一组由有限个向量组成的基的线性空间称为有限维的；否则称为无限维的。）

例如，\(\mathbb{R}^{n}\)的维数是\(n\)。向量组
\(\left( 1 , 0 , … , 0 \right) , \left( 0 , 1 , 0 , … , 0 \right) , … , \left( 0 , … , 0 , 1 \right)\)
线性无关且张成整个\(\mathbb{R}^{n}\)，这两点都不难验证，因此它是\(\mathbb{R}^{n}\)的一组基（称为标准基）。全体平面向量组成的空间（按坐标即\(\mathbb{R}^{2}\)）的维数是\(2\)，全体空间向量组成的空间（按坐标即\(\mathbb{R}^{3}\)）的维数是\(3\)。平时所说的“二维”“三维”，就是按维数来说的。一般地，维数就是坐标的个数：在\(n\)维线性空间中，每个向量都要用\(n\)个实数来度量。

并不是所有线性空间都是有限维的。例如在函数空间中，函数\(1 , x , x^{2} , … , x^{m}\)无论\(m\)多大都线性无关：一个非零的多项式不可能有无穷多个根，所以若它们的线性组合恒等于零函数，组合系数就只能全为零。于是这个空间含有任意多个线性无关的向量，再由上面的事实（线性无关的向量不可能多于张成它的向量组），它不可能由有限个向量张成。也就是说，函数空间是无限维的。无限维空间的例子还有很多，这里不再列举。

有了基与维数，线性空间中的计算大都可以化为坐标的计算：取定一组基，每个向量有唯一的坐标，向量的运算化为坐标的运算。下一节讨论线性空间之间的映射，即线性变换。那里会看到，一个线性变换完全由它把一组基中的向量映到何处所决定。

\begin{exercise}
（1）判断下列\(\mathbb{R}^{3}\)的子集是否为子空间：\\
（a）\(\left\{ \left( x , y , z \right) \mid x + y + z = 0 \right\}\)；\\
（b）\(\left\{ \left( x , y , z \right) \mid x + y + z = 1 \right\}\)；\\
（c）\(\left\{ \left( x , y , z \right) \mid x y = 0 \right\}\)；\\
（d）\(\left\{ \left( x , y , z \right) \mid x = 2 y , z = 3 y \right\}\)。

（2）判断下列向量组是否线性相关：（a）\(\left( 1 , 2 \right)\)，\(\left( 2 , 4 \right)\)；（b）\(\left( 1 , 0 \right)\)，\(\left( 0 , 1 \right)\)；（c）\(\left( 1 , 1 , 0 \right)\)，\(\left( 0 , 1 , 1 \right)\)，\(\left( 1 , 0 , 1 \right)\)。

（3）在\(\mathbb{R}^{2}\)中，验证\(e_{1} = \left( 1 , 2 \right)\)与\(e_{2} = \left( 1 , 3 \right)\)构成一组基，并求\(\left( 4 , 7 \right)\)在这组基下的坐标。

（4）判断向量\(\left( 17 , - 4 , 5 \right)\)是否属于\(\left( 2 , 1 , - 3 \right)\)与\(\left( 1 , - 2 , 4 \right)\)张成的子空间。

（5）设\(a , b , c\)线性无关，证明\(a + b\)，\(b + c\)，\(c + a\)也线性无关。
\end{exercise}

\section{线性变换*}
所谓线性变换，其关键理解在于什么是线性。变换是好理解的，即把一个东西映射成另一个东西，比如把一个向量变成另一个向量，把一个空间变成另一个空间。而所谓线性，即是要保证变换的可加性与齐次性。总的来说，可以定义变换\(T\)为线性变换，当且仅当\(T\)满足
\begin{equation}
    T \left( \alpha x + \beta y \right) := \alpha T \left( x \right) + \beta T \left( y \right)
\end{equation}
其中\(\alpha , \beta\)为标量，\(x , y \in V\)。容易看出，其和“直线”并无太大关联，而只要求映射保证与向量的加法、数乘“可交换”。进一步地，如果考虑直线\(f \left( x \right) = k x\)，可知\(f \left( \alpha x \right) = k \left( \alpha x \right) = \alpha \left( k x \right) = \alpha f \left( x \right)\)（可加性\(f \left( x + y \right) = f \left( x \right) + f \left( y \right)\)同样显然成立），从而\(f\)是一个线性变换，将\(x\)变换为\(k x\)。进一步地，考虑二维空间中的两个变换：
\begin{equation}
    T := \left( x , y \right) \mapsto \left( 2 x , 3 y \right)
\end{equation}
\begin{equation}
    T := \left( x , y \right) \mapsto \left( x + y , 2 y \right)
\end{equation}
显然，它们都是将二维空间变换成另一个二维空间的线性变换（一般称之为\(\mathbb{R}^{2}\)到\(\mathbb{R}^{2}\)的线性变换）。其中第一个便是伸缩变换，如果有一个圆，那么变换之后它就变成了椭圆；第二个则可以将一个长方形拉成平行四边形。

如果考虑\(\left( x , y \right)\)为二维空间中的一个向量，自然地\(\left( 2 x , 3 y \right) , \left( x + y , 2 y \right)\)也都是二维空间中的向量。因此这个变换将一个二维空间中的所有向量映射为另一组向量（这样的映射称为线性映射）。那么参考\(y = k x\)将\(x\)变换为\(y\)，我们能否找到一种运算模式，它可以通过\(y = K x\)来描述这样的变换呢？

\begin{theorem}\label{thm:linmap-basis}
设\(e_{1} , … , e_{n}\)是线性空间\(V\)的一组基，则\(V\)上的每个线性变换\(T\)都由像\(T \left( e_{1} \right) , … , T \left( e_{n} \right)\)完全决定；反过来，\(V\)中任意\(n\)个向量都可以作为一个线性变换在\(e_{1} , … , e_{n}\)处的像。
\end{theorem}

证明是直接的。对\(V\)中任意向量\(x = x_{1} e_{1} + \cdots + x_{n} e_{n}\)，由线性性有
\(T \left( x \right) = T \left( x_{1} e_{1} + \cdots + x_{n} e_{n} \right) = x_{1} T \left( e_{1} \right) + \cdots + x_{n} T \left( e_{n} \right)\)，
所以只要知道\(T\)把一组基中的向量映到何处，\(T\)就完全确定了。反过来，任意给定\(n\)个向量\(w_{1} , … , w_{n}\)，按
\(T \left( x_{1} e_{1} + \cdots + x_{n} e_{n} \right) = x_{1} w_{1} + \cdots + x_{n} w_{n}\)
定义的映射是线性变换，并且以它们为像。也就是说，\(V\)上的线性变换与一组基的\(n\)个像一一对应。每个像向量有\(n\)个坐标，合起来是\(n^{2}\)个数（这里是二维向量，取\(n = 2\)，所以是\(4\)个数）。那么容易定义二维情况下
\begin{equation}
    K = \left[ \begin{matrix}k_{11} & k_{12} \\ k_{21} & k_{22}\end{matrix} \right]
\end{equation}
并满足如下运算规则：
\begin{equation}
    y = K x = \left[ \begin{matrix}k_{11} & k_{12} \\ k_{21} & k_{22}\end{matrix} \right] \left[ \begin{matrix}x_{1} \\ x_{2}\end{matrix} \right] = \left[ \begin{matrix}k_{11} x_{1} + k_{12} x_{2} \\ k_{21} x_{1} + k_{22} x_{2}\end{matrix} \right] = \left[ \begin{matrix}y_{1} \\ y_{2}\end{matrix} \right]
\end{equation}
这样，便满足了任意形式的线性变换的运算。这样的情形可以自然而然地推广到高维向量，而\(K\)就是所谓的线性变换矩阵。矩阵的第\(j\)列，就是第\(j\)个基向量像的坐标；把各像的坐标按列排起来，就得到矩阵。

以下介绍5种二维空间中的线性变换。应当注意，在高维空间中仍然存在类似作用或性质的变换，这里不详细讲解，有兴趣可以阅读《Linear Algebra Done Right》作为线性代数的入门。

\paragraph{1°伸缩变换} \mbox{}

在之前的学习中，已经知道对坐标轴的伸缩是\(x^{'} = \lambda_{x} x , y^{'} = \lambda_{y} y\)。那么容易知道，此时的变换矩阵为\(\left[ \begin{matrix}\lambda_{x} & 0 \\ 0 & \lambda_{y}\end{matrix} \right]\)。

\paragraph{2°旋转变换} \mbox{}

如果将一个空间旋转\(\theta\)角，其变换矩阵为\(\left[ \begin{matrix}\cos \theta & - \sin \theta \\ \sin \theta & \cos \theta\end{matrix} \right]\)。该矩阵的推导过程主要是将向量的旋转变为标准正交基的旋转，进而计算两组系数。

\paragraph{3°翻折变换} \mbox{}

翻折变换指关于\(y = x\)对称，即交换\(x , y\)坐标，那么容易知道变换矩阵为\(\left[ \begin{matrix}0 & 1 \\ 1 & 0\end{matrix} \right]\)。

\paragraph{4°投影变换} \mbox{}

把向量投影在\(\theta\)角的直线上，仍然是采用旋转变换中类似的推导方式，考虑一组标准正交基的基向量的投影，容易知道该变换矩阵为\(\left[ \begin{matrix}\cos^{2} \theta & \cos \theta \sin \theta \\ \cos \theta \sin \theta & \sin^{2} \theta\end{matrix} \right]\)。

\paragraph{5°镜像变换} \mbox{}

镜像变换是更一般的翻折变换，即对称轴为任意过原点的与\(x\)轴成\(\theta\)角的直线。其变换矩阵为\(\left[ \begin{matrix}2 \cos^{2} \theta - 1 & 2 \cos \theta \sin \theta \\ 2 \cos \theta \sin \theta & 2 \sin^{2} \theta - 1\end{matrix} \right]\)。

大体上，高中数学中常用的线性变换为这五种，而在几何题中比较常见的一般有伸缩、旋转、投影，需要用到翻折的较少，需要用到镜像的很少见。

五种变换里已经出现过核的例子：投影变换把垂直于那条直线的向量全部映成零向量，这些向量组成一条过原点的直线。一般地，使\(T \left( x \right) = \vec{0}\)的向量组成的集合称为\(T\)的核，它是\(V\)的一个子空间；前面齐次线性方程组的解集就是它的一个例子。

{\kaiti{}思考\quad 平移\(T \left( x \right) = x + b\)（\(b \ne \vec{0}\)）为什么不是线性变换？}

\begin{example}
（1）设\(T\)是\(\mathbb{R}^{2}\)上的线性变换，已知\(T \left( 1 , 0 \right) = \left( 1 , 2 \right)\)，\(T \left( 0 , 1 \right) = \left( 3 , - 1 \right)\)，求\(T \left( 2 , 1 \right)\)与\(T\)的矩阵。

（2）求投影到\(x\)轴的变换的核，并说明它的几何意义。

（3）单位圆\(x^{2} + y^{2} = 1\)在伸缩变换\(\left( x , y \right) \mapsto \left( 2 x , y \right)\)下变成了什么曲线？
\end{example}

\begin{exercise}
（1）判断下列映射是否为\(\mathbb{R}^{2}\)上的线性变换：（a）\(T \left( x , y \right) = \left( x + 1 , y \right)\)；（b）\(T \left( x , y \right) = \left( x^{2} , y \right)\)；（c）\(T \left( x , y \right) = \left( 2 x - y , x + 3 y \right)\)。

（2）写出下列变换的矩阵：（a）把横坐标拉伸为原来的\(3\)倍，纵坐标不变；（b）绕原点旋转90°；（c）投影到\(x\)轴上。

（3）求投影到直线\(y = x\)上的变换的核，并说明理由。

（4）设\(T\)是线性变换，\(T \left( e_{1} \right) = \left( 2 , - 1 \right)\)，\(T \left( e_{2} \right) = \left( 1 , 1 \right)\)，求\(T \left( 3 , 2 \right)\)与\(T\)的矩阵。

（5）证明：线性变换把过原点的直线仍映成过原点的直线（也可能映成原点）。
\end{exercise}
\newpage